Table \ref{tab:results} shows the performance of all models on four datasets. It also shows the relative improvements of our model compared to the strongest baselines.
% Among the baselines, MARank performs the worst, presumably because it only looks at the most recent items of the user. TiSASRec outperforms SASRec, due to its usage of time intervals which makes its attention operation more comprehensive. As reported in \cite{BERT4Rec}, BERT4Rec outperforms SASRec due to its use of bi-directional transformers (as opposed to uni-directional transformers), which gives stronger prediction capabilities. BERT4Rec even outperforms TiSASRec in almost all cases. This is interesting considering that BERT4Rec does not utilize time interval information at all.
Our model gives the best performance in all metrics and datasets. On average, MEANTIME achieves \textbf{9.07\%}, \textbf{6.36\%}, \textbf{10.43\%}, and \textbf{9.14\%} improvements over the strongest baselines in Recall@5, Recall@10, NDCG@5 and NDCG@10, respectively. Given that we tune all models under the same set of hyperparameters (including the number of attention heads), we can deduce that the performance boost comes from the novel multi-temporal embeddings and attention operations that we proposed. Especially, although TiSASRec uses additional temporal information, it does not show impressive improvement over SASRec and BERT4Rec. We argue that using a single embedding scheme throughout all attention heads was a hindrance to their potential improvement. In contrast, each attention head takes different temporal embedding to extract specialized patterns in our design. This joint design of diverse temporal embeddings and attention mechanism was the key to our improvement.
